\section{Problem Statement: The Bridge Responsibility Gap}
\label{sec:problem}
% ================================================================================

\begin{definition}[Bridge component]\label{def:bridge}
A \emph{bridge} is any component $B$ that mediates between a client
transport $T_c$ and an agent runtime $A$, such that messages from
$T_c$ reach $A$ only via $B$ and messages from $A$ reach $T_c$ only
via $B$. The bridge is the locus of identity binding (it attaches
\texttt{Agent-ID} and \texttt{Authority-Scope} from AGTP to agent
messages), of audit emission (it emits events that a GAR consumer
observes), and of session-state continuity (it records the resume point
that PACE-style re-attach starts from).
\end{definition}

\begin{definition}[Bridge Responsibility Gap]\label{def:gap}
The \emph{Bridge Responsibility Gap} is the \emph{phenomenon} of
orphaned obligations at the composition boundary of AGTP, GAR, and PACE.
It is concretely realised as the \emph{set} of obligations comprising
gap classes G1--G4, defined below.

\begin{description}[leftmargin=2em]
\item[G1. Client-Detach-But-Agent-Continues.] The client transport $T_c$
disconnects (mobile handoff, network failure, user-initiated background)
while the agent $A$ is mid-execution. The bridge must record a
resume point, suspend the forwarding of new tool calls that arrive
while no client is attached (refusing them, holding them for an answer
after re-attach, or forwarding only those the operator allowlisted,
Section~\ref{sec:state-machine}), and leave work that is already in
flight undisturbed.

\item[G2. Human-Escalation-Pending.] The agent escalates a decision to a
human in the loop (HITL), but the human is unavailable (sleep, time zone,
context switch). The bridge must set the human dimension of the session to pending, so that all
subsequent tool calls are queued rather than executed, and
must emit a \texttt{HUMAN\_PENDING} governance event. HEM already
specifies a pending state in which every state transition is rejected, with
timeouts and an escalation chain (Section~\ref{sec:rw-hem}); what G2 adds
is the queue of tool calls held while the client may also be away, and the
relation of the human wait to a detach.

\item[G3. Resume-After-Detach.] The client re-attaches
(PACE-style) after a detach. The bridge must authenticate the client
by a challenge it answers under a registered key, verify that the
history the client presents is a prefix of the bridge's chain, must transition to the re-attaching ($\textsc{Res}$)
attachment, must re-verify the hash chain over governance events, and
must transition back to attached only if verification succeeds.

\item[G4. Bridge-Restart-With-Live-Sessions.] The bridge process
restarts (crash, upgrade) while sessions are live. Every client
transport is lost, but the agent and the human dimension are not. The
bridge must rebuild each session from its durable event log, treat it as
detached, and record a resume point, without losing or duplicating
governance events or queued tool calls.
\end{description}
\end{definition}

The Bridge Responsibility Gap is not a hypothetical concern. AgentThread
\cite{zheng2026formal} reports that ``no protocol assigns
enforcement for cross-protocol behavior'' and that ``[i]nsecurity in
agent protocols is therefore not only a specification or implementation
problem, but also a responsibility gap across protocols, SDKs, and
deployments'' (it also finds that ``composition is
universally orphaned''). AgentThread's analysis constructs
cases in which each action appears tool-valid in isolation while the
harmful intent is visible only in the composed trace. It analyses other
agent protocols --- MCP, A2A, ANP, and two different protocols both
abbreviated ``ACP'' in the source paper (the Agent Communication
Protocol and the Agent Client Protocol; see Section~\ref{sec:rw-agtp-thread}) --- and does not analyse AGTP, GAR or PACE, so its finding is
motivation here, not a result about them. The gap classes G1--G4 above are
this document's own derivation of the obligations that none of AGTP, GAR and PACE, nor GAR's companion HEM, as read, assigns to a bridge; the derivation is not shown to be
complete. The Bridge Contract is a specification that owns them. The
``Why this obligation is out of <draft>'s scope'' paragraphs in
Section~\ref{sec:rw-agtp}--\ref{sec:rw-pace} and the Agent Audit
Trail subsection argue that, in the versions read, the gap is one of scope (a scope
mismatch) rather than an oversight: a revision of any cited draft that
owned these obligations would be expanding its scope, not closing a gap
inside it. This is an argument from a reading of dated texts, not a
prediction about those drafts, and it weakens if one of them is
extended to own these duties (Section~\ref{sec:limitations}). It is the
reason given for the Bridge Contract being a separate specification, not
a profile of an existing draft.

\begin{example}[Concrete scenario]
A mobile user asks an agent to reconcile three months of corporate
expenses. The agent enqueues forty tool calls (database queries,
spreadsheet operations, email reads). After twenty minutes the user's
phone moves between cell towers; the transport disconnects for forty
minutes. (This scenario uses the \emph{persistent} profile,
$T_{\text{detach}} = 7$~days; under the default
\emph{interactive-short} profile, $T_{\text{detach}} = 30$~minutes,
the session would terminate before the user returns.)
During those forty minutes the agent finishes the calls that
were already in flight, escalates one
ambiguity to the user (which the user does not see for forty minutes),
has the calls that depend on the answer queued by the bridge,
and attempts to extend its authority scope to read an additional mailbox
``to save a round trip''. Without a Bridge Contract, all of these
behaviours are unspecified: the runtime may silently extend authority
(a violation of the obligation that I3 captures), may execute the
escalated tool call (a violation of the obligation that I2 captures),
or may lose the partial state produced during the detach (a consequence
of violating I1's emission requirement). With a conformant bridge,
each of these behaviours is audited and recoverable, and each is
refused when it passes through the bridge; a change the runtime makes
on its own, without asking the bridge, is observed and reported but not
prevented (Section~\ref{sec:coverage}).
This scenario exercises transitions
$(\textsc{Att},\textsc{None}) \to (\textsc{Det},\textsc{None}) \to
(\textsc{Det},\textsc{Pend}) \to (\textsc{Res},\textsc{Pend}) \to
(\textsc{Att},\textsc{Pend})$ (see Figure~\ref{fig:state-machine}; see
Table~\ref{tab:transitions} for preconditions and postconditions).
\end{example}

\paragraph{Decision guide: when the Bridge Contract helps and when it
is overhead.}
The Bridge Contract adds value only when at least one of the gap
classes G1--G4 can fire in the deployment. Concretely, it helps when
\emph{any} of these holds:
\begin{itemize}[leftmargin=2em]
\item \textbf{Long-lived agent + mobile client.} The agent outlives
the interactive client that initiated it (G1, G3): a mobile user, a
browser tab that can close, a laptop that can sleep. The bridge
keeps the session governable across the detach.
\item \textbf{Human-in-the-loop with an unavailable human.} The agent
escalates to a human who may be asleep, in another time zone, or
context-switched away (G2). The bridge queues, times out, and
replays on resolution.
\item \textbf{Multi-device principal.} The principal has several
devices attached at once or in turn (I15). The bridge treats the
session as detached only when none is connected, and a second device
is a join, not a resume.
\item \textbf{Regulated audit without a GEC.} The deployment needs a
tamper-evident audit trail but has no Governing Enforcement Component
(Definition~\ref{def:gecless}). The bridge is the event source, in
GAR-shaped form.
\item \textbf{Bridge restart with live sessions.} The bridge process
may crash or upgrade while sessions are live (G4). The bridge
rebuilds from the durable log and treats every session as detached.
\end{itemize}
The contract is \emph{overhead}---every state transition emits an
event, every detach produces a resume record---when \emph{all} of
these hold:
\begin{itemize}[leftmargin=2em]
\item short-lived sessions (the agent finishes within one client
connection);
\item always-connected client (no detach is possible);
\item always-available human (no escalation ever waits);
\item single-device principal (no client set needed);
\item GEC-present deployment (the GEC is the audit source);
\item no bridge restarts in scope (single-shot execution).
\end{itemize}
Conformant deployments MAY skip resume-record production when no
detach is possible, but MUST still emit lifecycle events
(\texttt{ATTACHED}, \texttt{TERMINATION}). The decision is a
deployment-level configuration, not a property of the contract.

\subsection{Scope and Failure Model}\label{sec:scope}

The Bridge Contract specifies the obligations of a bridge component
during the gap classes G1--G4 of Definition~\ref{def:gap}. The
following failure modes and compositions are \emph{in scope}:

\begin{itemize}[leftmargin=2em]
\item Client transport disconnect (full disconnect, both directions
down), with the agent continuing to execute. (G1)
\item Human-in-the-loop escalation with the human currently unavailable
(sleep, time zone, context switch). (G2)
\item Client re-attach after a prior disconnect, with verification of
the presented history and the hash chain. (G3)
\item Bridge restart with live sessions, recovered from a durable event
log. (G4)
\item Authentication of the re-attaching client, including replay of captured re-attach requests and attempts to
terminate a session with forged ones. (G3)
\item Several clients of one principal, attached at once or in turn:
the session is detached only when none is connected (client set,
Section~\ref{sec:product-model}).
\item Clean session termination from any state.
\item Authority extension attempts in any cell: always rejected when
requested through the bridge, and reported as \emph{observed} when the
runtime announces the change itself (I3).
\item A \emph{Byzantine client}: one that replays captured
re-attach requests, forges signatures, claims a history that is not a
prefix of the chain, or tries to terminate a session with forged
re-attach attempts, provided no registered private key is compromised
(A7). Such attempts are rejected without a state change (I4, I9, I11),
with one stated exception: a caller that holds the live-nonce limit can keep a detached client from
re-attaching until the detach timeout ends the session (limitations
register, Table~\ref{tab:limitations}).
\end{itemize}

The following failure modes and compositions are \emph{out of scope}
for this document and are deferred to future revisions or to other
specifications:

\begin{itemize}[leftmargin=2em]
\item \textbf{Bridge availability.} Replication, failover and
storage faults beyond Axiom A1 are deployment-level concerns. The
contract covers restart of a single bridge with a durable log
(Section~\ref{sec:recovery}), not high availability.
\item \textbf{Sharing between principals.} A session belongs to one
principal. Several principals attached to the same session, with
per-principal authority, are not modelled.
\item \textbf{Network partition (partial detach).} The contract models
full disconnect of a client only; partial connectivity (e.g., one direction down)
is out of scope.
\item \textbf{Concurrent events.} The state machine assumes serialised
event emission per session; concurrent emission of, e.g.,
\texttt{HUMAN\_PENDING} and \texttt{DETACHMENT} is not modelled as
true concurrency. Because the product machine has orthogonal dimensions
(Section~\ref{sec:product-model}), serialised events from different
dimensions commute on the cell they reach (Axiom A5).
\item \textbf{Compromised keys.} Theft of a client's private key (Axiom A7). A Byzantine client without the key
is in scope (above); a client that holds the key is indistinguishable
from the legitimate one.
\item \textbf{Byzantine bridge.} A bridge that deliberately suppresses
governance events or forges events is out of scope.
Mitigations require external transparency logs (e.g., SCITT) and are
deployment-level.
\item \textbf{Multi-hop composition (BC $\circ$ BC).} If two bridges
are composed in sequence, the obligations at their composition boundary
are not modelled here. Multi-hop composition is the target for a future
revision.
\item \textbf{Concurrent BC and GEC.} In a deployment with both a
bridge and a GEC, the GEC remains the sole authoritative event source
(per GAR \S 3); the bridge forwards events to the GEC for signing.
The interaction is governed by the GEC-less deployment profile
(Definition~\ref{def:gecless}) and Invariants I-SIGN-EVENT and I-SIGN-JWKS
(Section~\ref{sec:invariants-signing}).
\item \textbf{Multi-agent composition (A2A, MCP).} Declared out of
scope; A2A and MCP operate above the session layer governed by the
Bridge Contract.
\item \textbf{QUIC 0-RTT replay.} AGTP-SESSION \S 10.5 and the
replay-safety profile of AGTP-BINDINGS \S 4.6 are authoritative for replay
protection of the underlying transport; the
Bridge Contract defers to them and adds Invariant I8
(Section~\ref{sec:invariants-quic}) for the bridge-layer consequence.
\end{itemize}

A future revision may extend the contract to cover partial detach,
true concurrency, replication, and multi-hop composition.

\subsection{What the bridge can see: coverage}\label{sec:coverage}

The obligations of this contract attach to what crosses the bridge.
Two consequences are part of the specification, not caveats to it.
\begin{itemize}[leftmargin=2em]
\item \textbf{Calls submitted for approval.} I2 and I13 apply to the
tool calls that the runtime \emph{submits to the bridge} for approval or
forwarding. In some runtimes the bridge sees only the permission request
(for example the Agent Client Protocol's
\texttt{session/request\_permission}; throughout this document ``ACP''
means the Agent Client Protocol, see Section~\ref{sec:rw-agtp-thread}), not the
execution that follows. Execution that the bridge does not mediate is
outside the contract: the bridge cannot queue, suspend or audit what it
does not see.
\item \textbf{Enforced and observed authority changes.} A scope or mode
change \emph{requested through the bridge} is enforced: widening is
refused (I3). A change that the runtime announces itself, as an Agent
Client Protocol (ACP) agent announces its own mode change, is
\emph{observed}: the bridge
cannot prevent it, so it records \texttt{AUTHORITY\_EXTENSION\_REJECTED}
(source \texttt{observed}) and SHOULD instruct the runtime to revert.
Between the announcement and the revert the runtime holds authority the
bridge has not granted; the contract bounds this by audit and by the
bridge's own record of the scope, which never widens, not by prevention.
A bridge claims \emph{enforced} authority only where the runtime routes
all authority changes through it; where the runtime may change authority
on its own, the bridge's authority is \emph{observed}, not enforced.
The deployment documents which case holds.
\end{itemize}

\subsection{Formal status of this specification}\label{sec:formal-status}

This document is a \emph{normative engineering specification with a
bounded formal model}. It provides
definitions, invariants, axioms, lemmas, and a safety argument, and
it accompanies them with a TLA+ model that is checked on bounded
configurations. The wire format (Appendix~\ref{app:wire-format}) is
normative for client--bridge messages in the signed bytes and the reason
names, which the vectors and the known-answer files check; the vectors
compare committed event kinds and final cells, not response bytes, so
passing V01--V137 does not make two implementations byte-compatible at the
message level (Section~\ref{sec:limitations}). The unit of
conformance is the committed event log and the final session state:
two implementations that produce the same committed event kinds and
the same final cell on the same input trace are conformant at the
semantic level, whether or not they share a transport binding. The
normative state model is the product
machine of Section~\ref{sec:product-model}; the names ATTACHED, DETACHED, PENDING\_HUMAN, RESUMING and TERMINATED
survive only as a derived, lossy view
(Section~\ref{sec:mece-status}). A TLA+ model of the product machine
accompanies the document (Appendix~\ref{app:formal}); TLC finds no
invariant violation on bounded configurations (two sessions with two
principals and two permissions, and a one-session deep run), and a set of
deliberately broken variants of the model is rejected by TLC. The result
does not extend to unbounded sessions (Appendix~\ref{app:formal-limitations}).

% ================================================================================
